\chapter{Diffusion Models: Intuition}

\section{Noise Injection Processes}

Diffusion models are generative models based on gradually transforming data into noise and then reversing this process.

\medskip

\noindent
\textbf{Forward process.}  
We define a sequence of random variables $(x_0, x_1, \dots, x_T)$ where:
\begin{itemize}
    \item $x_0 \sim p_{\text{data}}$
    \item $x_T$ approaches a simple distribution (e.g., Gaussian noise)
\end{itemize}

\medskip

\noindent
\textbf{Noise injection.}  
At each step, we add small Gaussian noise:
\[
x_t = \sqrt{1 - \beta_t} \, x_{t-1} + \sqrt{\beta_t} \, \varepsilon_t,
\quad \varepsilon_t \sim \mathcal{N}(0, I).
\]

\medskip

\noindent
\textbf{Markov structure.}
\begin{itemize}
    \item Each step depends only on the previous state
    \item The process is a Markov chain
\end{itemize}

\medskip

\noindent
\textbf{Effect.}
\begin{itemize}
    \item Gradually destroys structure in data
    \item Transforms data distribution into noise
\end{itemize}

\medskip

\noindent
\textbf{Insight.}  
The forward process defines a controlled way to map complex data distributions to simple ones.

\section{Denoising Perspective}

The generative process learns to reverse the noise injection.

\medskip

\noindent
\textbf{Reverse process.}  
We aim to model:
\[
p_\theta(x_{t-1} \mid x_t),
\]
which removes noise step by step.

\medskip

\noindent
\textbf{Denoising interpretation.}
\begin{itemize}
    \item Each step predicts a less noisy version of the data
    \item The model learns to recover structure from noise
\end{itemize}

\medskip

\noindent
\textbf{Learning objective.}  
The model is trained to predict the noise added at each step:
\[
\varepsilon_\theta(x_t, t) \approx \varepsilon_t.
\]

\medskip

\noindent
\textbf{Simplified loss.}
\[
\mathbb{E}_{t, x_0, \varepsilon}
\left[
\| \varepsilon - \varepsilon_\theta(x_t, t) \|^2
\right].
\]

\medskip

\noindent
\textbf{Sampling.}
\begin{itemize}
    \item Start from noise $x_T \sim \mathcal{N}(0, I)$
    \item Iteratively denoise to obtain $x_0$
\end{itemize}

\medskip

\noindent
\textbf{Insight.}  
Generation is achieved by repeatedly denoising random noise into structured data.

\section{Score-Based View}

Diffusion models can also be interpreted through score functions.

\medskip

\noindent
\textbf{Score function.}
\[
s(x) = \nabla_x \log p(x).
\]

\medskip

\noindent
\textbf{Key idea.}  
Instead of modeling $p(x)$ directly, learn its score.

\medskip

\noindent
\textbf{Connection to diffusion.}
\begin{itemize}
    \item The denoising function approximates the score of noisy distributions
    \item Each step corresponds to estimating $\nabla_x \log p_t(x)$
\end{itemize}

\medskip

\noindent
\textbf{Langevin dynamics.}
\[
x_{t+1} = x_t + \eta \, \nabla_x \log p(x_t) + \sqrt{2\eta} \, \xi_t,
\]
where $\xi_t$ is Gaussian noise.

\medskip

\noindent
\textbf{Interpretation.}
\begin{itemize}
    \item Samples are generated by following gradients of log-density
    \item Noise ensures exploration of the distribution
\end{itemize}

\medskip

\noindent
\textbf{Insight.}  
Diffusion models connect denoising with gradient-based sampling from probability distributions.

\section{Conceptual Connections}

Diffusion models unify several ideas from probability, optimization, and physics.

\medskip

\noindent
\textbf{Connections to other models.}
\begin{itemize}
    \item \textbf{Latent variable models:} introduce intermediate noisy states
    \item \textbf{Energy-based models:} relate to score functions
    \item \textbf{Markov processes:} sequential transformations
\end{itemize}

\medskip

\noindent
\textbf{Continuous-time limit.}
\begin{itemize}
    \item Discrete steps approximate stochastic differential equations
    \item Leads to a continuous diffusion process
\end{itemize}

\medskip

\noindent
\textbf{Geometry of data.}
\begin{itemize}
    \item Noise smooths the data distribution
    \item Denoising recovers structure
\end{itemize}

\medskip

\noindent
\textbf{Comparison with other generative models.}
\begin{itemize}
    \item VAEs: probabilistic latent modeling
    \item GANs: adversarial training
    \item Flows: invertible transformations
    \item Diffusion: iterative refinement
\end{itemize}

\medskip

\noindent
\textbf{Insight.}  
Diffusion models combine probabilistic modeling with iterative denoising and gradient-based sampling.

\section{A Unifying Perspective}

Diffusion models can be understood through multiple lenses:

\begin{itemize}
    \item \textbf{Probabilistic:} modeling distributions through noise processes
    \item \textbf{Algorithmic:} iterative denoising procedures
    \item \textbf{Geometric:} following gradients in data space
\end{itemize}

\medskip

\noindent
\textbf{Core components.}
\begin{itemize}
    \item Forward noise process
    \item Reverse denoising process
    \item Learned score or noise predictor
\end{itemize}

\medskip

\noindent
\textbf{Key insight.}  
Complex data distributions can be generated by reversing a simple stochastic process.

\section{Outlook}

This chapter introduced diffusion models from an intuitive perspective:
\begin{itemize}
    \item noise injection processes,
    \item denoising as generation,
    \item score-based interpretations,
    \item conceptual connections to other models.
\end{itemize}

\medskip

In the next chapter, we study diffusion models more rigorously, including their probabilistic formulation and training objectives.

\medskip

\noindent
\textbf{Guiding principle.}  
Complex structure can be learned by progressively refining simple noise through learned transformations.